\chapter*{前言与资源使用说明}
\addcontentsline{toc}{chapter}{前言与资源使用说明}

% SOURCE_PDF_PAGE 5

\hypertarget{source-p005}{}
司t
\sourceblock{80.23}{images-full/p005-d000.png}
全国大学生数学竞赛是由中国数学会主办、面向本科学生的全国性高水平学科竞赛，
旨在激励大学生学习数学的兴趣，培养他们分析问题、解决问题的能力，为青年学子搭建
一个展示数学思维能力和学习成果的平台.
全国大学生数学竞赛自2009 年开展以来，得到了全国各赛区和各高校的大力支持和
帮助.在各赛区和承办单位的辛勤努力下，该赛事己经连续成功举办了15 届.参赛高校
由第一届的400 多所增加到第十五届的1100 余所，参赛人数也由第一届的2 万多人增
加到第十五届的近29 万人，在全国高校产生了广泛的影响.这15 届的承办单位分别是
国防科技大学、北京航壁航天大学、同济大学、电子科技大学、中国科学技术大学、华中
科技大学、福建师范大学、北京科技大学、西安交通大学、哈尔滨工业大学、武汉大学、
吉林大学、华东师范大学、广东工业大学、山东大学.在此，向全国各赛区、各高校和各
\sourceblock{230.52}{images-full/p005-d001.png}
全国大学生数学竞赛分初赛和决赛两个阶段，每个阶段都包含数学专业和非数学专
业两大类别.随着全国大学生数学竞赛持续深入开展，参赛学生越来越多，规模越来越大，
参赛学生和竞赛指导教师对数学竞赛资料的需求也越来越大，但是目前专门针对大学生
数学竞赛的辅导教材不多，满足不了参赛学生和竞赛指导教师日益增长的需求.为了帮助
大学生和热爱数学的人士更好地了解这项全国性赛事，并对有志参加数学竞赛的大学生
进行专业的数学竞赛指导，经过与科学出版社协商，全国大学生数学竞赛命题组决定撰写
《全国大学生数学竞赛解析教程(数学专业类) >>和《全国大学生数学竞赛解析教程(非数
学专业类)>> ，作为《全国大学生数学竞赛参赛指南》的拓展参考读物.
《全国大学生数学竞赛解析教程(数学专业类) >>紧扣《全国大学生数学竞赛参赛指
南》里的全国大学生数学竞赛数学专业类考试内容，分上、下两册出版，上册是数学分析
的相关内容，下册是高等代数和解析几何的相关内容.数学分析部分共七章，包括Euclid
空间，极限与连续，微分，级数，Riemann 积分、曲线积分及曲面积分，反常积分及含参变
量积分，综合与拓展.高等代数部分共三章，包括矩阵与线性变换、线性方程组与二次型、
多项式与矩阵方程.解析几何部分共三章，包括向量代数、几何空间中的平面与直线、二
次曲面与二次曲线.这些内容都是经过命题组专家精心组织的，充分体现了大学生数学竞
赛试题的综合性、高阶性、创新性、挑战性等特点.
《全国大学生数学竞赛解析教程(数学专业类) >>具有如下几方面的特色:
\par\noindent (1) 强调基础知识.对于竞赛中必考的核心内容与知识点，特别是基本概念、基本理


% SOURCE_PDF_PAGE 6

\hypertarget{source-p006}{}
论和基本方法，在每个章节都首先作简要总结，然后对一些精选的例题和赛题进行系统详
细的阐述和深入浅出的讲解，有些例题还给出了必要的注.
\sourceblock{432.22}{images-full/p006-d000.png}
易懂的典型方法，期望读者能够举一反三.
\par\noindent (3) 注重综合技能.大学生数学竞赛不仅测试相对容易的试题，而且更多的问题都具
有一定的灵活性和挑战性，需要具有较强的综合技能才能正确求解.因此，书中各部分内
容均体现出对综合技能的足够重视.
\par\noindent (4) 揭示数学思想.在全国大学生数学竞赛中，许多试题不仅技巧性强、富于原创性，
而且其求解方法也往往彰显出诸多奇思妙想，因此在写法上尽力去揭示隐藏在技巧背后
的数学思想.这有利于读者学习和应用数学方法，提高数学思维能力，更是培养科学思维
品质的"利器"、激励科学创新的源泉.
《全国大学生数学竞赛解析教程(数学专业类)>>可作为高等院校数学专业类大学生参
加全国大学生数学竞赛的备考辅导教程，也可作为这些学生提开解题能力的进阶读物，还
可作为广大考研学子的复习资料.
限于时间，不当之处在所难免，敬请读者提出宝贵意见.
全国大学生数学竞赛命题组
2023 年12 月11 日
..


% SOURCE_PDF_PAGE 7

\hypertarget{source-p007}{}
资源使用说明
1.使用流程
\par\noindent (1) 刮开封底激活码的涂层，使用微信扫描二维码. 根据提示，注册登录到中科助学
通平台，激活本书的配套资源.
\sourceblock{450.00}{images-full/p007-d000.png}
2. 内容介绍
为了让读者对全国大学生数学竞赛有更多更生动的理解，读者购买正版教材之后，可
获取我们提供的两类配套资源，一类是视频资源，另一类是文本资源.
\par\noindent (1) 视频资源来源于全国大学生数学竞赛命题组专家们关于竞赛的专题讲解报告，本
书包含11 个视频，内容包括矩阵、线性变换、解析几何等.
\par\noindent (2) 文本资源来源于竞赛模拟测试题，本书包含22 道题，可用于模拟训练，并提供详
细参考解答，供备考使用.
两类资源总的资源二维码如下. 另外，涉及各章节知识的资源二维码也会放在相应的
章后.
\sourceblock{190.97}{images-full/p007-d001.png}
